% Part: first-order-logic
% Chapter: proof-systems
% Section: tableaux

\documentclass[../../../include/open-logic-section]{subfiles}

\begin{document}

\iftag{FOL}
      {\olfileid{fol}{prf}{tab}}
      {\olfileid{pl}{prf}{tab}}

\olsection{\usetoken{P}{tableau}}

అనేక \tetoken{వ్యుత్పత్తి}{!!{derivation}} వ్యవస్థలు \tetoken{వాక్యాలు}{!!{sentence}s} అమరికలతో పనిచేస్తే,
\tetoken{టాబ్లోలు}{!!{tableau}s} అనేవి \tetoken{చిహ్నిత సూత్రాలు}{!!{signed formula}s} అనే చిహ్నిత సూత్రాలతో పనిచేస్తాయి.
\tetoken{ఒక చిహ్నిత సూత్రం}{!!^a{signed formula}} అంటే ఒక సత్యమూల్య సంకేతం ($\True$ లేదా
$\False$), మరియు \tetoken{ఒక వాక్యంతో}{!!a{sentence}} కూడిన జత:
\[
\sFmla{\True}{!A} \text{ or } \sFmla{\False}{!A}.
\]
కిందివైపుకు శాఖలుగా విడిపోయే వృక్షంలో అమర్చిన \tetoken{చిహ్నిత సూత్రాలతో}{!!{signed formula}s}
ఒక \tetoken{టాబ్లో}{!!^a{tableau}} ఏర్పడుతుంది. అది కొన్ని \emph{పరికల్పనలతో}
మొదలై, వాటి పైనున్న \tetoken{చిహ్నిత సూత్రాలలో}{!!{signed formula}s} ఒకదానికి నిగమన నియమాల్లో
ఒకదాన్ని వర్తింపజేయడం వల్ల వచ్చే \tetoken{చిహ్నిత సూత్రాలతో}{!!{signed formula}s} కొనసాగుతుంది.
ప్రతి నియమం ఒక శాఖ చివరికి ఒకటి లేదా అంతకంటే ఎక్కువ \tetoken{చిహ్నిత సూత్రాలను}{!!{signed
  formula}s} చేర్చనిస్తుంది, లేదా రెండు \tetoken{చిహ్నిత సూత్రాలను}{!!{signed formula}s} పక్కపక్కనే
చేర్చనిస్తుంది. రెండవ సందర్భంలో శాఖ రెండుగా చీలి, కొత్తగా చేర్చిన
రెండు \tetoken{చిహ్నిత సూత్రాలు}{!!{signed formula}s} ఆ రెండు శాఖల చివరలుగా ఏర్పడతాయి.

ఒక సంక్లిష్ట \tetoken{చిహ్నిత సూత్రానికి}{!!{signed formula}} నియమాన్ని వర్తింపజేస్తే, దాని
తక్షణ ఉప-\tetoken{సూత్రాలు}{!!{formula}s} అయిన \tetoken{చిహ్నిత సూత్రాలను}{!!{signed formula}s} చేర్చాలి. ప్రతి
రెండు సంకేతాల్లో ఒక్కొక్కదానికి ఒక నియమం చొప్పున అవి జతలుగా ఉంటాయి.
ఉదాహరణకు, $\TRule{\True}{\land}$ నియమం
$\sFmla{\True}{!A \land !B}$కు వర్తిస్తుంది;
$\sFmla{\True}{!A \land !B}$ను కలిగి ఉన్న ఏ
శాఖ చివరికైనా $\sFmla{\True}{!A}$, $\sFmla{\True}{!B}$ అనే రెండు
\tetoken{చిహ్నిత సూత్రాలనూ}{!!{signed formula}s} చేర్చనిస్తుంది. $\TRule{\False}{\land}$
నియమం $\sFmla{\False}{!A}$, $\sFmla{\False}{!B}$లను పక్కపక్కనే
చేర్చి ఒక శాఖను రెండుగా చీల్చనిస్తుంది.
\sourcecorrection{OLTEPRF-002}{ఆంగ్ల మూలం $\TRule{\False}{\land}$
నియమం రెండవ ఆర్గ్యుమెంటుగా సంయోగ సంకేతం బదులు పూర్తి
$!A \land !B$ సూత్రాన్ని ఇచ్చింది. స్థూలనిర్దేశంలోని
\texttt{TRule\{Sign\}\{Op\}} ఒప్పందం, పక్కనే ఉన్న సత్య-సంయోగ
నియమం ప్రకారం ఆర్గ్యుమెంటును $\land$గా సరిచేశాం.}
ప్రతి శాఖలో $\sFmla{\True}{!A}$, $\sFmla{\False}{!A}$ అనే సరిపోలే
\tetoken{చిహ్నిత సూత్రాలు}{!!{signed formula}s} జత ఉంటే ఒక \tetoken{టాబ్లో}{!!^a{tableau}} సంవృతమైనది.

\tetoken{టాబ్లోలు}{!!{tableau}s} ఆధారంగా $\Proves$ సంబంధాన్ని ఇలా నిర్వచిస్తాం:
$\Gamma_0 = \{!B_1, \dots, !B_n\} \subseteq \Gamma$ అనే ఏదో
పరిమిత సమితి ఉండి, కింది పరికల్పనలకు ఒక సంవృత \tetoken{టాబ్లో}{!!{tableau}} ఉంటే,
అప్పుడు మాత్రమే $\Gamma \Proves !A$:
\[
\{\sFmla{\False}{!A}, \sFmla{\True}{!B_1}, \dots, \sFmla{\True}{!B_n}\}
\]
ఉదాహరణకు, $\Proves (!A \land !B) \lif !A$ అని చూపే సంవృత
\tetoken{టాబ్లో}{!!{tableau}} ఇది:
\begin{oltableau}
  [\sFmla{\False}{(\formula{A} \land \formula{B}) \lif \formula{A}}, just = \TAss
    [\sFmla{\True}{\formula{A} \land \formula{B}}, just = {\TRule{\False}{\lif}[1]}
      [\sFmla{\False}{\formula{A}}, just = {\TRule{\False}{\lif}[1]}
        [\sFmla{\True}{\formula{A}}, just = {\TRule{\True}{\land}[2]}
          [\sFmla{\True}{\formula{B}}, just = {\TRule{\True}{\land}[2]}, close]
        ]
      ]
    ]
  ]
\end{oltableau}
\sourcecorrection{OLTEPRF-003}{ఆంగ్ల మూలంలోని ఉదాహరణ, రెండవ వరుసలోని
సత్య-సంయోగాన్ని $\sFmla{\True}{!A}$, $\sFmla{\True}{!B}$లుగా
విడదీస్తూనే, రెండు ఫలితాలకూ సత్య-షరతీయ నియమపు గుర్తు ఇచ్చింది.
ఫలితాలను నియంత్రించే పూర్వ సూత్రం ప్రకారం రెండు గుర్తులనూ
సత్య-సంయోగ నియమంగా సరిచేసి, వరుస సంఖ్య~$2$ను అలాగే ఉంచాం.}

ప్రతి సూచికకు $!B_i \in \Gamma$ అయి, కింది పరికల్పనలకు ఒక సంవృత
\tetoken{టాబ్లో}{!!{tableau}} ఉంటే, సమితి $\Gamma$ \tetoken{టాబ్లో}{!!{tableau}} కలనంలో
వైరుధ్యభరితమైనది:
\[
\{\sFmla{\True}{!B_1}, \dots, \sFmla{\True}{!B_n}\}
\]
\sourcecorrection{OLTEPRF-004}{ప్రదర్శించిన పరికల్పనలు అన్నీ
$\Gamma$ నుంచి రావాల్సి ఉండగా, ఆంగ్ల మూలం “ఏదో ఒక $!B_i$” మాత్రమే
$\Gamma$లో ఉండాలని చెప్పింది. ముందున్న నిగమ్యత నిర్వచనంలోని
పరిమిత ఉపసమితి షరతుకు అనుగుణంగా, ప్రతి సూచికకు సభ్యత్వం అవసరమని
సరిచేశాం.}

1950లలో ఎవర్ట్ బెత్, యాకో హింటిక్కా ఒకరికొకరు స్వతంత్రంగా
\tetoken{టాబ్లోలను}{!!^{tableau}s} కనుగొన్నారు; రేమండ్ స్మల్యన్ వాటిని సరళీకరించి
ప్రాచుర్యంలోకి తెచ్చాడు. ఒక \tetoken{టాబ్లోను}{!!a{tableau}} నిర్మించడం ఎంతో
క్రమబద్ధమైన ప్రక్రియ కాబట్టి వాటిని ఉపయోగించడం చాలా సులభం.
\tetoken{టాబ్లోలు}{!!{tableau}s} క్రమబద్ధ స్వభావం వల్ల వాటిని కంప్యూటరులో అమలు చేయడం
కూడా అనుకూలం. అయితే, ఒక \tetoken{టాబ్లోను}{!!a{tableau}} చదవడం తరచుగా కష్టం;
నిరూపణలతో దాని సంబంధం కొన్నిసార్లు సులభంగా కనిపించదు. ఈ విధానం
ఎంతో సాధారణమైనది కూడా; అనేక వేర్వేరు తర్కవ్యవస్థలకు \tetoken{టాబ్లో}{!!{tableau}}
వ్యవస్థలు ఉన్నాయి. ఇచ్చిన \tetoken{వాక్యాలు}{!!{sentence}s} (సమితులను) సంతృప్తిపరిచే
\tetoken{నిర్మాణాలను}{!!{structure}s} కనుగొనడంలోనూ \tetoken{టాబ్లోలు}{!!^{tableau}s} సహాయపడతాయి: సమితి
సంతృప్తిపరచదగినదైతే దానికి సంవృత \tetoken{టాబ్లో}{!!{tableau}} ఉండదు; అంటే ఏ
\tetoken{టాబ్లోలోనైనా}{!!{tableau}} ఒక వివృత శాఖ ఉంటుంది. ఒక శాఖపై వర్తింపజేయగల
ప్రతి నియమాన్నీ దానిపై వర్తింపజేసి ఉంటే, ఆ వివృత శాఖ నుంచి
సంతృప్తిపరిచే \tetoken{నిర్మాణాన్ని}{!!{structure}} “చదివి తీసుకోవచ్చు.” సీక్వెంట్
కలనంతో కూడా ఎంతో సన్నిహిత సంబంధం ఉంది: సారాంశంగా, సంవృత
\tetoken{టాబ్లో}{!!{tableau}} అంటే తలకిందులుగా రాసిన, సీక్వెంట్ కలనంలోని సంక్షిప్త
\tetoken{వ్యుత్పత్తి}{!!{derivation}}.

\end{document}
